
\documentclass{article}
\usepackage{colortbl}
\usepackage{makecell}
\usepackage{multirow}
\usepackage{supertabular}

\begin{document}

\newcounter{utterance}

\centering \large Interaction Transcript for game `codenames', experiment `opponent\_difficult', episode 3 with gemma{-}4{-}26b{-}a4b{-}it.
\vspace{24pt}

{ \footnotesize  \setcounter{utterance}{1}
\setlength{\tabcolsep}{0pt}
\begin{supertabular}{c@{$\;$}|p{.15\linewidth}@{}p{.15\linewidth}p{.15\linewidth}p{.15\linewidth}p{.15\linewidth}p{.15\linewidth}}
    \# & $\;$A & \multicolumn{4}{c}{Game Master} & $\;\:$B\\
    \hline

    \theutterance \stepcounter{utterance}  
    & & & \multicolumn{2}{p{0.3\linewidth}}{
        \cellcolor[rgb]{0.9,0.9,0.9}{
            \makecell[{{p{\linewidth}}}]{
                \texttt{\tiny{[GM$|$GM]}}
                {'hidden': {'team': ['акула', 'триъгълник', 'ученик', 'шип', 'токио', 'бермуди', 'жак', 'време', 'преса'], 'innocent': ['танц', 'игла', 'промяна', 'мечка', 'пролет', 'низ', 'обувка'], 'opponent': ['сърце', 'отрова', 'лед', 'кетчуп', 'болт', 'кенгуру', 'планина', 'кръгъл'], 'assassin': ['лекар']}, 'revealed': {'team': {'team': [], 'innocent': [], 'opponent': [], 'assassin': []}, 'opponent': {'team': [], 'innocent': [], 'opponent': [], 'assassin': []}}}
            }
        }
    }
    & & \\ \\

    \theutterance \stepcounter{utterance}  
    & & \multicolumn{4}{p{0.6\linewidth}}{
        \cellcolor[rgb]{0.9,0.9,0.9}{
            \makecell[{{p{\linewidth}}}]{
                \texttt{\tiny{[P1$\langle$GM]}}
                \texttt{Дайте една подсказка с една дума, която е свързана с някои от думите на вашия екип от списъка по{-}долу. Можете да избирате с колко и с кои думи да е свързана подсказката ви.} \\
\texttt{Целта е партньорът ви да отгатне всички думи на вашия отбор въз основа на подсказките, които му давате, преди другият отбор да е разкрил всички свои думи. Ако партньорът ви отгатне дума на "противника" или "разсейваща" дума, неговият ход за отгатване приключва и противниковият отбор получава възможност да направи ход. Ако партньорът ви отгатне дума "убиец", играта приключва и отборът ви губи незабавно.} \\
\texttt{Колкото повече думи свързва вашата подсказка, толкова по{-}бързо можете да спечелите играта, но асоциациите може да се окажат по{-}трудни за отгатване от партньора ви, така че избирайте подсказващата дума разумно. Подсказващата дума трябва да е семантично свързана с целевите думи, не може да е една от думите в списъците или да съдържа части от тях.} \\
\\ 
\texttt{Винаги давайте подсказката от една дума и списъка със свързаните с нея целеви думи, разделени със запетая, в следния формат и правете отговорите си възможно най{-}кратки, като никога не включвате друг текст, освен този, който се изисква в тази форма:} \\
\\ 
\texttt{ПОДСКАЗКА: <ДУМА>} \\
\texttt{ЦЕЛИ: <ЦЕЛИ>} \\
\\ 
\texttt{Думите на вашия отбор са: акула, триъгълник, ученик, шип, токио, бермуди, жак, време, преса.} \\
\texttt{Думите на отбора на опонента ви са: сърце, отрова, лед, кетчуп, болт, кенгуру, планина, кръгъл.} \\
\texttt{Разсейващите думи са: танц, игла, промяна, мечка, пролет, низ, обувка.} \\
\texttt{Думите убийци са: лекар.} \\
            }
        }
    }
    & \\ \\

    \theutterance \stepcounter{utterance}  
    & \multicolumn{4}{p{0.6\linewidth}}{
        \cellcolor[rgb]{0.9,0.9,0.9}{
            \makecell[{{p{\linewidth}}}]{
                \texttt{\tiny{[P1$\rangle$GM]}}
                \texttt{ПОДСКАЗКА: Остри} \\
\texttt{ЦЕЛИ: акула, триъгълник, шип, жак} \\
            }
        }
    }
    & & \\ \\

    \theutterance \stepcounter{utterance}  
    & & & \multicolumn{2}{p{0.3\linewidth}}{
        \cellcolor[rgb]{0.9,0.9,0.9}{
            \makecell[{{p{\linewidth}}}]{
                \texttt{\tiny{[GM$|$GM]}}
                \texttt{остри} \\
            }
        }
    }
    & & \\ \\

    \theutterance \stepcounter{utterance}  
    & & & \multicolumn{2}{p{0.3\linewidth}}{
        \cellcolor[rgb]{0.9,0.9,0.9}{
            \makecell[{{p{\linewidth}}}]{
                \texttt{\tiny{[GM$|$GM]}}
                ['акула', 'триъгълник', 'шип', 'жак']
            }
        }
    }
    & & \\ \\

    \theutterance \stepcounter{utterance}  
    & & & \multicolumn{2}{p{0.3\linewidth}}{
        \cellcolor[rgb]{0.9,0.9,0.9}{
            \makecell[{{p{\linewidth}}}]{
                \texttt{\tiny{[GM$|$GM]}}
                {'word': 'акула', 'assignment': 'team'}
            }
        }
    }
    & & \\ \\

    \theutterance \stepcounter{utterance}  
    & & & \multicolumn{2}{p{0.3\linewidth}}{
        \cellcolor[rgb]{0.9,0.9,0.9}{
            \makecell[{{p{\linewidth}}}]{
                \texttt{\tiny{[GM$|$GM]}}
                {'word': 'триъгълник', 'assignment': 'team'}
            }
        }
    }
    & & \\ \\

    \theutterance \stepcounter{utterance}  
    & & & \multicolumn{2}{p{0.3\linewidth}}{
        \cellcolor[rgb]{0.9,0.9,0.9}{
            \makecell[{{p{\linewidth}}}]{
                \texttt{\tiny{[GM$|$GM]}}
                {'word': 'шип', 'assignment': 'team'}
            }
        }
    }
    & & \\ \\

    \theutterance \stepcounter{utterance}  
    & & & \multicolumn{2}{p{0.3\linewidth}}{
        \cellcolor[rgb]{0.9,0.9,0.9}{
            \makecell[{{p{\linewidth}}}]{
                \texttt{\tiny{[GM$|$GM]}}
                {'word': 'жак', 'assignment': 'team'}
            }
        }
    }
    & & \\ \\

    \theutterance \stepcounter{utterance}  
    & & \multicolumn{4}{p{0.6\linewidth}}{
        \cellcolor[rgb]{0.9,0.9,0.9}{
            \makecell[{{p{\linewidth}}}]{
                \texttt{\tiny{[GM$\rangle$P2]}}
                \texttt{Представете списък, разделен със запетаи, от до 4 думи от следния списък, които най{-}добре се отнасят или са най{-}тясно свързани с думата 'остри'. Винаги започвайте списъка си с предположения с ОТГОВОР: и не включвайте друг текст в отговора си.} \\
\\ 
\texttt{промяна, шип, бермуди, токио, танц, сърце, лед, кръгъл, кетчуп, акула, планина, ученик, обувка, отрова, лекар, низ, кенгуру, време, триъгълник, жак, мечка, игла, пролет, болт, преса} \\
            }
        }
    }
    & \\ \\

    \theutterance \stepcounter{utterance}  
    & & & \multicolumn{4}{p{0.6\linewidth}}{
        \cellcolor[rgb]{0.9,0.9,0.9}{
            \makecell[{{p{\linewidth}}}]{
                \texttt{\tiny{[GM$\langle$P2]}}
                \texttt{ОТГОВОР: шип, игла, акула, триъгълник} \\
            }
        }
    }
     \\ \\

    \theutterance \stepcounter{utterance}  
    & & & \multicolumn{2}{p{0.3\linewidth}}{
        \cellcolor[rgb]{0.9,0.9,0.9}{
            \makecell[{{p{\linewidth}}}]{
                \texttt{\tiny{[GM$|$GM]}}
                ['шип', 'игла', 'акула', 'триъгълник']
            }
        }
    }
    & & \\ \\

    \theutterance \stepcounter{utterance}  
    & & & \multicolumn{2}{p{0.3\linewidth}}{
        \cellcolor[rgb]{0.9,0.9,0.9}{
            \makecell[{{p{\linewidth}}}]{
                \texttt{\tiny{[GM$|$GM]}}
                {'word': 'шип', 'assignment': 'team'}
            }
        }
    }
    & & \\ \\

    \theutterance \stepcounter{utterance}  
    & & & \multicolumn{2}{p{0.3\linewidth}}{
        \cellcolor[rgb]{0.9,0.9,0.9}{
            \makecell[{{p{\linewidth}}}]{
                \texttt{\tiny{[GM$|$GM]}}
                {'word': 'шип', 'assignment': 'team'}
            }
        }
    }
    & & \\ \\

    \theutterance \stepcounter{utterance}  
    & & & \multicolumn{2}{p{0.3\linewidth}}{
        \cellcolor[rgb]{0.9,0.9,0.9}{
            \makecell[{{p{\linewidth}}}]{
                \texttt{\tiny{[GM$|$GM]}}
                {'word': 'игла', 'assignment': 'innocent'}
            }
        }
    }
    & & \\ \\

    \theutterance \stepcounter{utterance}  
    & & & \multicolumn{2}{p{0.3\linewidth}}{
        \cellcolor[rgb]{0.9,0.9,0.9}{
            \makecell[{{p{\linewidth}}}]{
                \texttt{\tiny{[GM$|$GM]}}
                \texttt{игла} \\
            }
        }
    }
    & & \\ \\

    \theutterance \stepcounter{utterance}  
    & & & \multicolumn{2}{p{0.3\linewidth}}{
        \cellcolor[rgb]{0.9,0.9,0.9}{
            \makecell[{{p{\linewidth}}}]{
                \texttt{\tiny{[GM$|$GM]}}
                {'word': 'лед', 'assignment': 'opponent'}
            }
        }
    }
    & & \\ \\

    \theutterance \stepcounter{utterance}  
    & & & \multicolumn{2}{p{0.3\linewidth}}{
        \cellcolor[rgb]{0.9,0.9,0.9}{
            \makecell[{{p{\linewidth}}}]{
                \texttt{\tiny{[GM$|$GM]}}
                {'word': 'кенгуру', 'assignment': 'opponent'}
            }
        }
    }
    & & \\ \\

    \theutterance \stepcounter{utterance}  
    & & & \multicolumn{2}{p{0.3\linewidth}}{
        \cellcolor[rgb]{0.9,0.9,0.9}{
            \makecell[{{p{\linewidth}}}]{
                \texttt{\tiny{[GM$|$GM]}}
                {'hidden': {'team': ['акула', 'триъгълник', 'ученик', 'токио', 'бермуди', 'жак', 'време', 'преса'], 'innocent': ['танц', 'промяна', 'мечка', 'пролет', 'низ', 'обувка'], 'opponent': ['сърце', 'отрова', 'кетчуп', 'болт', 'планина', 'кръгъл'], 'assassin': ['лекар']}, 'revealed': {'team': {'team': ['шип'], 'innocent': ['игла'], 'opponent': [], 'assassin': []}, 'opponent': {'team': [], 'innocent': [], 'opponent': ['лед', 'кенгуру'], 'assassin': []}}}
            }
        }
    }
    & & \\ \\

    \theutterance \stepcounter{utterance}  
    & & \multicolumn{4}{p{0.6\linewidth}}{
        \cellcolor[rgb]{0.9,0.9,0.9}{
            \makecell[{{p{\linewidth}}}]{
                \texttt{\tiny{[P1$\langle$GM]}}
                \texttt{Думата шип беше позната правилно. Думата игла беше позната, но е неутрален дума. Ходът на вашия отборник завърши там.} \\
\texttt{Сега дайте още една подсказка, свързана с някои от останалите думи на вашия отбор, и списък на свързаните целеви думи. Не забравяйте да започнете подсказката си с 'ПОДСКАЗКА: ', поставете нов ред и започнете списъка с целевите думи, разделен със запетая, с 'ЦЕЛИ: '. Забележете: някои думи са премахнати от списъците в сравнение с предишните заявки.} \\
\\ 
\texttt{Останалите думи на вашия отбор са: акула, триъгълник, ученик, токио, бермуди, жак, време, преса.} \\
\texttt{Останалите думи на вашия противник са: сърце, отрова, лед, кетчуп, болт, кенгуру, планина, кръгъл.} \\
\texttt{Останалите отвличащи думи са: танц, промяна, мечка, пролет, низ, обувка.} \\
\texttt{Останалите убийствени думи са: лекар.} \\
            }
        }
    }
    & \\ \\

    \theutterance \stepcounter{utterance}  
    & \multicolumn{4}{p{0.6\linewidth}}{
        \cellcolor[rgb]{0.9,0.9,0.9}{
            \makecell[{{p{\linewidth}}}]{
                \texttt{\tiny{[P1$\rangle$GM]}}
                \texttt{ПОДСКАЗКА: Япония} \\
\texttt{ЦЕЛИ: токио, акула, бермуди} \\
            }
        }
    }
    & & \\ \\

    \theutterance \stepcounter{utterance}  
    & & & \multicolumn{2}{p{0.3\linewidth}}{
        \cellcolor[rgb]{0.9,0.9,0.9}{
            \makecell[{{p{\linewidth}}}]{
                \texttt{\tiny{[GM$|$GM]}}
                \texttt{япония} \\
            }
        }
    }
    & & \\ \\

    \theutterance \stepcounter{utterance}  
    & & & \multicolumn{2}{p{0.3\linewidth}}{
        \cellcolor[rgb]{0.9,0.9,0.9}{
            \makecell[{{p{\linewidth}}}]{
                \texttt{\tiny{[GM$|$GM]}}
                ['токио', 'акула', 'бермуди']
            }
        }
    }
    & & \\ \\

    \theutterance \stepcounter{utterance}  
    & & & \multicolumn{2}{p{0.3\linewidth}}{
        \cellcolor[rgb]{0.9,0.9,0.9}{
            \makecell[{{p{\linewidth}}}]{
                \texttt{\tiny{[GM$|$GM]}}
                {'word': 'токио', 'assignment': 'team'}
            }
        }
    }
    & & \\ \\

    \theutterance \stepcounter{utterance}  
    & & & \multicolumn{2}{p{0.3\linewidth}}{
        \cellcolor[rgb]{0.9,0.9,0.9}{
            \makecell[{{p{\linewidth}}}]{
                \texttt{\tiny{[GM$|$GM]}}
                {'word': 'акула', 'assignment': 'team'}
            }
        }
    }
    & & \\ \\

    \theutterance \stepcounter{utterance}  
    & & & \multicolumn{2}{p{0.3\linewidth}}{
        \cellcolor[rgb]{0.9,0.9,0.9}{
            \makecell[{{p{\linewidth}}}]{
                \texttt{\tiny{[GM$|$GM]}}
                {'word': 'бермуди', 'assignment': 'team'}
            }
        }
    }
    & & \\ \\

    \theutterance \stepcounter{utterance}  
    & & \multicolumn{4}{p{0.6\linewidth}}{
        \cellcolor[rgb]{0.9,0.9,0.9}{
            \makecell[{{p{\linewidth}}}]{
                \texttt{\tiny{[GM$\rangle$P2]}}
                \texttt{Думата шип беше позната правилно. Думата игла беше позната, но е неутрален дума. Вашият ход завърши там.} \\
\texttt{Сега предоставете още един списък, разделен със запетаи, от най{-}малко 1 и до 3 думи от следния списък с думи, които най{-}добре се отнасят или са най{-}тясно свързани с думата 'япония'. Не забравяйте да започнете отговора си с ОТГОВОР:. Забележете: някои думи са премахнати от списъка в сравнение с предишните заявки.} \\
\\ 
\texttt{промяна, бермуди, токио, танц, сърце, кръгъл, кетчуп, акула, планина, ученик, обувка, отрова, лекар, низ, време, триъгълник, жак, мечка, пролет, болт, преса} \\
            }
        }
    }
    & \\ \\

    \theutterance \stepcounter{utterance}  
    & & & \multicolumn{4}{p{0.6\linewidth}}{
        \cellcolor[rgb]{0.9,0.9,0.9}{
            \makecell[{{p{\linewidth}}}]{
                \texttt{\tiny{[GM$\langle$P2]}}
                \texttt{ОТГОВОР: токио} \\
            }
        }
    }
     \\ \\

    \theutterance \stepcounter{utterance}  
    & & & \multicolumn{2}{p{0.3\linewidth}}{
        \cellcolor[rgb]{0.9,0.9,0.9}{
            \makecell[{{p{\linewidth}}}]{
                \texttt{\tiny{[GM$|$GM]}}
                ['токио']
            }
        }
    }
    & & \\ \\

    \theutterance \stepcounter{utterance}  
    & & & \multicolumn{2}{p{0.3\linewidth}}{
        \cellcolor[rgb]{0.9,0.9,0.9}{
            \makecell[{{p{\linewidth}}}]{
                \texttt{\tiny{[GM$|$GM]}}
                {'word': 'токио', 'assignment': 'team'}
            }
        }
    }
    & & \\ \\

    \theutterance \stepcounter{utterance}  
    & & & \multicolumn{2}{p{0.3\linewidth}}{
        \cellcolor[rgb]{0.9,0.9,0.9}{
            \makecell[{{p{\linewidth}}}]{
                \texttt{\tiny{[GM$|$GM]}}
                {'word': 'токио', 'assignment': 'team'}
            }
        }
    }
    & & \\ \\

    \theutterance \stepcounter{utterance}  
    & & & \multicolumn{2}{p{0.3\linewidth}}{
        \cellcolor[rgb]{0.9,0.9,0.9}{
            \makecell[{{p{\linewidth}}}]{
                \texttt{\tiny{[GM$|$GM]}}
                {'word': 'кетчуп', 'assignment': 'opponent'}
            }
        }
    }
    & & \\ \\

    \theutterance \stepcounter{utterance}  
    & & & \multicolumn{2}{p{0.3\linewidth}}{
        \cellcolor[rgb]{0.9,0.9,0.9}{
            \makecell[{{p{\linewidth}}}]{
                \texttt{\tiny{[GM$|$GM]}}
                {'word': 'планина', 'assignment': 'opponent'}
            }
        }
    }
    & & \\ \\

    \theutterance \stepcounter{utterance}  
    & & & \multicolumn{2}{p{0.3\linewidth}}{
        \cellcolor[rgb]{0.9,0.9,0.9}{
            \makecell[{{p{\linewidth}}}]{
                \texttt{\tiny{[GM$|$GM]}}
                {'hidden': {'team': ['акула', 'триъгълник', 'ученик', 'бермуди', 'жак', 'време', 'преса'], 'innocent': ['танц', 'промяна', 'мечка', 'пролет', 'низ', 'обувка'], 'opponent': ['сърце', 'отрова', 'болт', 'кръгъл'], 'assassin': ['лекар']}, 'revealed': {'team': {'team': ['шип', 'токио'], 'innocent': ['игла'], 'opponent': [], 'assassin': []}, 'opponent': {'team': [], 'innocent': [], 'opponent': ['лед', 'кенгуру', 'кетчуп', 'планина'], 'assassin': []}}}
            }
        }
    }
    & & \\ \\

    \theutterance \stepcounter{utterance}  
    & & \multicolumn{4}{p{0.6\linewidth}}{
        \cellcolor[rgb]{0.9,0.9,0.9}{
            \makecell[{{p{\linewidth}}}]{
                \texttt{\tiny{[P1$\langle$GM]}}
                \texttt{Думата токио беше позната правилно. Ходът на вашия отборник завърши там.} \\
\texttt{Сега дайте още една подсказка, свързана с някои от останалите думи на вашия отбор, и списък на свързаните целеви думи. Не забравяйте да започнете подсказката си с 'ПОДСКАЗКА: ', поставете нов ред и започнете списъка с целевите думи, разделен със запетая, с 'ЦЕЛИ: '. Забележете: някои думи са премахнати от списъците в сравнение с предишните заявки.} \\
\\ 
\texttt{Останалите думи на вашия отбор са: акула, триъгълник, ученик, бермуди, жак, време, преса.} \\
\texttt{Останалите думи на вашия противник са: сърце, отрова, кетчуп, болт, планина, кръгъл.} \\
\texttt{Останалите отвличащи думи са: танц, промяна, мечка, пролет, низ, обувка.} \\
\texttt{Останалите убийствени думи са: лекар.} \\
            }
        }
    }
    & \\ \\

    \theutterance \stepcounter{utterance}  
    & \multicolumn{4}{p{0.6\linewidth}}{
        \cellcolor[rgb]{0.9,0.9,0.9}{
            \makecell[{{p{\linewidth}}}]{
                \texttt{\tiny{[P1$\rangle$GM]}}
                \texttt{ПОДСКАЗКА: Мистерия} \\
\texttt{ЦЕЛИ: бермуди, акула, триъгълник} \\
            }
        }
    }
    & & \\ \\

    \theutterance \stepcounter{utterance}  
    & & & \multicolumn{2}{p{0.3\linewidth}}{
        \cellcolor[rgb]{0.9,0.9,0.9}{
            \makecell[{{p{\linewidth}}}]{
                \texttt{\tiny{[GM$|$GM]}}
                \texttt{мистерия} \\
            }
        }
    }
    & & \\ \\

    \theutterance \stepcounter{utterance}  
    & & & \multicolumn{2}{p{0.3\linewidth}}{
        \cellcolor[rgb]{0.9,0.9,0.9}{
            \makecell[{{p{\linewidth}}}]{
                \texttt{\tiny{[GM$|$GM]}}
                ['бермуди', 'акула', 'триъгълник']
            }
        }
    }
    & & \\ \\

    \theutterance \stepcounter{utterance}  
    & & & \multicolumn{2}{p{0.3\linewidth}}{
        \cellcolor[rgb]{0.9,0.9,0.9}{
            \makecell[{{p{\linewidth}}}]{
                \texttt{\tiny{[GM$|$GM]}}
                {'word': 'бермуди', 'assignment': 'team'}
            }
        }
    }
    & & \\ \\

    \theutterance \stepcounter{utterance}  
    & & & \multicolumn{2}{p{0.3\linewidth}}{
        \cellcolor[rgb]{0.9,0.9,0.9}{
            \makecell[{{p{\linewidth}}}]{
                \texttt{\tiny{[GM$|$GM]}}
                {'word': 'акула', 'assignment': 'team'}
            }
        }
    }
    & & \\ \\

    \theutterance \stepcounter{utterance}  
    & & & \multicolumn{2}{p{0.3\linewidth}}{
        \cellcolor[rgb]{0.9,0.9,0.9}{
            \makecell[{{p{\linewidth}}}]{
                \texttt{\tiny{[GM$|$GM]}}
                {'word': 'триъгълник', 'assignment': 'team'}
            }
        }
    }
    & & \\ \\

    \theutterance \stepcounter{utterance}  
    & & \multicolumn{4}{p{0.6\linewidth}}{
        \cellcolor[rgb]{0.9,0.9,0.9}{
            \makecell[{{p{\linewidth}}}]{
                \texttt{\tiny{[GM$\rangle$P2]}}
                \texttt{Думата токио беше позната правилно. Вашият ход завърши там.} \\
\texttt{Сега предоставете още един списък, разделен със запетаи, от най{-}малко 1 и до 3 думи от следния списък с думи, които най{-}добре се отнасят или са най{-}тясно свързани с думата 'мистерия'. Не забравяйте да започнете отговора си с ОТГОВОР:. Забележете: някои думи са премахнати от списъка в сравнение с предишните заявки.} \\
\\ 
\texttt{промяна, бермуди, танц, сърце, кръгъл, акула, ученик, обувка, отрова, лекар, низ, време, триъгълник, жак, мечка, пролет, болт, преса} \\
            }
        }
    }
    & \\ \\

    \theutterance \stepcounter{utterance}  
    & & & \multicolumn{4}{p{0.6\linewidth}}{
        \cellcolor[rgb]{0.9,0.9,0.9}{
            \makecell[{{p{\linewidth}}}]{
                \texttt{\tiny{[GM$\langle$P2]}}
                \texttt{ОТГОВОР: бермуди, отрова} \\
            }
        }
    }
     \\ \\

    \theutterance \stepcounter{utterance}  
    & & & \multicolumn{2}{p{0.3\linewidth}}{
        \cellcolor[rgb]{0.9,0.9,0.9}{
            \makecell[{{p{\linewidth}}}]{
                \texttt{\tiny{[GM$|$GM]}}
                ['бермуди', 'отрова']
            }
        }
    }
    & & \\ \\

    \theutterance \stepcounter{utterance}  
    & & & \multicolumn{2}{p{0.3\linewidth}}{
        \cellcolor[rgb]{0.9,0.9,0.9}{
            \makecell[{{p{\linewidth}}}]{
                \texttt{\tiny{[GM$|$GM]}}
                {'word': 'бермуди', 'assignment': 'team'}
            }
        }
    }
    & & \\ \\

    \theutterance \stepcounter{utterance}  
    & & & \multicolumn{2}{p{0.3\linewidth}}{
        \cellcolor[rgb]{0.9,0.9,0.9}{
            \makecell[{{p{\linewidth}}}]{
                \texttt{\tiny{[GM$|$GM]}}
                {'word': 'бермуди', 'assignment': 'team'}
            }
        }
    }
    & & \\ \\

    \theutterance \stepcounter{utterance}  
    & & & \multicolumn{2}{p{0.3\linewidth}}{
        \cellcolor[rgb]{0.9,0.9,0.9}{
            \makecell[{{p{\linewidth}}}]{
                \texttt{\tiny{[GM$|$GM]}}
                {'word': 'отрова', 'assignment': 'opponent'}
            }
        }
    }
    & & \\ \\

    \theutterance \stepcounter{utterance}  
    & & & \multicolumn{2}{p{0.3\linewidth}}{
        \cellcolor[rgb]{0.9,0.9,0.9}{
            \makecell[{{p{\linewidth}}}]{
                \texttt{\tiny{[GM$|$GM]}}
                \texttt{отрова} \\
            }
        }
    }
    & & \\ \\

    \theutterance \stepcounter{utterance}  
    & & & \multicolumn{2}{p{0.3\linewidth}}{
        \cellcolor[rgb]{0.9,0.9,0.9}{
            \makecell[{{p{\linewidth}}}]{
                \texttt{\tiny{[GM$|$GM]}}
                {'word': 'болт', 'assignment': 'opponent'}
            }
        }
    }
    & & \\ \\

    \theutterance \stepcounter{utterance}  
    & & & \multicolumn{2}{p{0.3\linewidth}}{
        \cellcolor[rgb]{0.9,0.9,0.9}{
            \makecell[{{p{\linewidth}}}]{
                \texttt{\tiny{[GM$|$GM]}}
                {'word': 'кръгъл', 'assignment': 'opponent'}
            }
        }
    }
    & & \\ \\

    \theutterance \stepcounter{utterance}  
    & & & \multicolumn{2}{p{0.3\linewidth}}{
        \cellcolor[rgb]{0.9,0.9,0.9}{
            \makecell[{{p{\linewidth}}}]{
                \texttt{\tiny{[GM$|$GM]}}
                {'hidden': {'team': ['акула', 'триъгълник', 'ученик', 'жак', 'време', 'преса'], 'innocent': ['танц', 'промяна', 'мечка', 'пролет', 'низ', 'обувка'], 'opponent': ['сърце'], 'assassin': ['лекар']}, 'revealed': {'team': {'team': ['шип', 'токио', 'бермуди'], 'innocent': ['игла'], 'opponent': ['отрова'], 'assassin': []}, 'opponent': {'team': [], 'innocent': [], 'opponent': ['лед', 'кенгуру', 'кетчуп', 'планина', 'болт', 'кръгъл'], 'assassin': []}}}
            }
        }
    }
    & & \\ \\

    \theutterance \stepcounter{utterance}  
    & & \multicolumn{4}{p{0.6\linewidth}}{
        \cellcolor[rgb]{0.9,0.9,0.9}{
            \makecell[{{p{\linewidth}}}]{
                \texttt{\tiny{[P1$\langle$GM]}}
                \texttt{Думата бермуди беше позната правилно. Думата отрова беше позната, но е противнически дума. Ходът на вашия отборник завърши там.} \\
\texttt{Сега дайте още една подсказка, свързана с някои от останалите думи на вашия отбор, и списък на свързаните целеви думи. Не забравяйте да започнете подсказката си с 'ПОДСКАЗКА: ', поставете нов ред и започнете списъка с целевите думи, разделен със запетая, с 'ЦЕЛИ: '. Забележете: някои думи са премахнати от списъците в сравнение с предишните заявки.} \\
\\ 
\texttt{Останалите думи на вашия отбор са: акула, триъгълник, ученик, жак, време, преса.} \\
\texttt{Останалите думи на вашия противник са: сърце, болт, кръгъл.} \\
\texttt{Останалите отвличащи думи са: танц, промяна, мечка, пролет, низ, обувка.} \\
\texttt{Останалите убийствени думи са: лекар.} \\
            }
        }
    }
    & \\ \\

    \theutterance \stepcounter{utterance}  
    & \multicolumn{4}{p{0.6\linewidth}}{
        \cellcolor[rgb]{0.9,0.9,0.9}{
            \makecell[{{p{\linewidth}}}]{
                \texttt{\tiny{[P1$\rangle$GM]}}
                \texttt{ПОДСКАЗКА: Геометрия} \\
\texttt{ЦЕЛИ: триъгълник, жак} \\
            }
        }
    }
    & & \\ \\

    \theutterance \stepcounter{utterance}  
    & & & \multicolumn{2}{p{0.3\linewidth}}{
        \cellcolor[rgb]{0.9,0.9,0.9}{
            \makecell[{{p{\linewidth}}}]{
                \texttt{\tiny{[GM$|$GM]}}
                \texttt{геометрия} \\
            }
        }
    }
    & & \\ \\

    \theutterance \stepcounter{utterance}  
    & & & \multicolumn{2}{p{0.3\linewidth}}{
        \cellcolor[rgb]{0.9,0.9,0.9}{
            \makecell[{{p{\linewidth}}}]{
                \texttt{\tiny{[GM$|$GM]}}
                ['триъгълник', 'жак']
            }
        }
    }
    & & \\ \\

    \theutterance \stepcounter{utterance}  
    & & & \multicolumn{2}{p{0.3\linewidth}}{
        \cellcolor[rgb]{0.9,0.9,0.9}{
            \makecell[{{p{\linewidth}}}]{
                \texttt{\tiny{[GM$|$GM]}}
                {'word': 'триъгълник', 'assignment': 'team'}
            }
        }
    }
    & & \\ \\

    \theutterance \stepcounter{utterance}  
    & & & \multicolumn{2}{p{0.3\linewidth}}{
        \cellcolor[rgb]{0.9,0.9,0.9}{
            \makecell[{{p{\linewidth}}}]{
                \texttt{\tiny{[GM$|$GM]}}
                {'word': 'жак', 'assignment': 'team'}
            }
        }
    }
    & & \\ \\

    \theutterance \stepcounter{utterance}  
    & & \multicolumn{4}{p{0.6\linewidth}}{
        \cellcolor[rgb]{0.9,0.9,0.9}{
            \makecell[{{p{\linewidth}}}]{
                \texttt{\tiny{[GM$\rangle$P2]}}
                \texttt{Думата бермуди беше позната правилно. Думата отрова беше позната, но е противнически дума. Вашият ход завърши там.} \\
\texttt{Сега предоставете още един списък, разделен със запетаи, от най{-}малко 1 и до 2 думи от следния списък с думи, които най{-}добре се отнасят или са най{-}тясно свързани с думата 'геометрия'. Не забравяйте да започнете отговора си с ОТГОВОР:. Забележете: някои думи са премахнати от списъка в сравнение с предишните заявки.} \\
\\ 
\texttt{промяна, танц, сърце, акула, ученик, обувка, лекар, низ, време, триъгълник, жак, мечка, пролет, преса} \\
            }
        }
    }
    & \\ \\

    \theutterance \stepcounter{utterance}  
    & & & \multicolumn{4}{p{0.6\linewidth}}{
        \cellcolor[rgb]{0.9,0.9,0.9}{
            \makecell[{{p{\linewidth}}}]{
                \texttt{\tiny{[GM$\langle$P2]}}
                \texttt{ОТГОВОР: триъгълник} \\
            }
        }
    }
     \\ \\

    \theutterance \stepcounter{utterance}  
    & & & \multicolumn{2}{p{0.3\linewidth}}{
        \cellcolor[rgb]{0.9,0.9,0.9}{
            \makecell[{{p{\linewidth}}}]{
                \texttt{\tiny{[GM$|$GM]}}
                ['триъгълник']
            }
        }
    }
    & & \\ \\

    \theutterance \stepcounter{utterance}  
    & & & \multicolumn{2}{p{0.3\linewidth}}{
        \cellcolor[rgb]{0.9,0.9,0.9}{
            \makecell[{{p{\linewidth}}}]{
                \texttt{\tiny{[GM$|$GM]}}
                {'word': 'триъгълник', 'assignment': 'team'}
            }
        }
    }
    & & \\ \\

    \theutterance \stepcounter{utterance}  
    & & & \multicolumn{2}{p{0.3\linewidth}}{
        \cellcolor[rgb]{0.9,0.9,0.9}{
            \makecell[{{p{\linewidth}}}]{
                \texttt{\tiny{[GM$|$GM]}}
                {'word': 'триъгълник', 'assignment': 'team'}
            }
        }
    }
    & & \\ \\

    \theutterance \stepcounter{utterance}  
    & & & \multicolumn{2}{p{0.3\linewidth}}{
        \cellcolor[rgb]{0.9,0.9,0.9}{
            \makecell[{{p{\linewidth}}}]{
                \texttt{\tiny{[GM$|$GM]}}
                {'word': 'сърце', 'assignment': 'opponent'}
            }
        }
    }
    & & \\ \\

    \theutterance \stepcounter{utterance}  
    & & & \multicolumn{2}{p{0.3\linewidth}}{
        \cellcolor[rgb]{0.9,0.9,0.9}{
            \makecell[{{p{\linewidth}}}]{
                \texttt{\tiny{[GM$|$GM]}}
                {'hidden': {'team': ['акула', 'ученик', 'жак', 'време', 'преса'], 'innocent': ['танц', 'промяна', 'мечка', 'пролет', 'низ', 'обувка'], 'opponent': [], 'assassin': ['лекар']}, 'revealed': {'team': {'team': ['шип', 'токио', 'бермуди', 'триъгълник'], 'innocent': ['игла'], 'opponent': ['отрова'], 'assassin': []}, 'opponent': {'team': [], 'innocent': [], 'opponent': ['лед', 'кенгуру', 'кетчуп', 'планина', 'болт', 'кръгъл', 'сърце'], 'assassin': []}}}
            }
        }
    }
    & & \\ \\

    \theutterance \stepcounter{utterance}  
    & & \multicolumn{4}{p{0.6\linewidth}}{
        \cellcolor[rgb]{0.9,0.9,0.9}{
            \makecell[{{p{\linewidth}}}]{
                \texttt{\tiny{[P1$\langle$GM]}}
                \texttt{Думата триъгълник беше позната правилно. Ходът на вашия отборник завърши там.} \\
\texttt{Сега дайте още една подсказка, свързана с някои от останалите думи на вашия отбор, и списък на свързаните целеви думи. Не забравяйте да започнете подсказката си с 'ПОДСКАЗКА: ', поставете нов ред и започнете списъка с целевите думи, разделен със запетая, с 'ЦЕЛИ: '. Забележете: някои думи са премахнати от списъците в сравнение с предишните заявки.} \\
\\ 
\texttt{Останалите думи на вашия отбор са: акула, ученик, жак, време, преса.} \\
\texttt{Останалите думи на вашия противник са: сърце.} \\
\texttt{Останалите отвличащи думи са: танц, промяна, мечка, пролет, низ, обувка.} \\
\texttt{Останалите убийствени думи са: лекар.} \\
            }
        }
    }
    & \\ \\

    \theutterance \stepcounter{utterance}  
    & \multicolumn{4}{p{0.6\linewidth}}{
        \cellcolor[rgb]{0.9,0.9,0.9}{
            \makecell[{{p{\linewidth}}}]{
                \texttt{\tiny{[P1$\rangle$GM]}}
                \texttt{ПОДСКАЗКА: Журналист} \\
\texttt{ЦЕЛИ: преса, ученик} \\
            }
        }
    }
    & & \\ \\

    \theutterance \stepcounter{utterance}  
    & & & \multicolumn{2}{p{0.3\linewidth}}{
        \cellcolor[rgb]{0.9,0.9,0.9}{
            \makecell[{{p{\linewidth}}}]{
                \texttt{\tiny{[GM$|$GM]}}
                \texttt{журналист} \\
            }
        }
    }
    & & \\ \\

    \theutterance \stepcounter{utterance}  
    & & & \multicolumn{2}{p{0.3\linewidth}}{
        \cellcolor[rgb]{0.9,0.9,0.9}{
            \makecell[{{p{\linewidth}}}]{
                \texttt{\tiny{[GM$|$GM]}}
                ['преса', 'ученик']
            }
        }
    }
    & & \\ \\

    \theutterance \stepcounter{utterance}  
    & & & \multicolumn{2}{p{0.3\linewidth}}{
        \cellcolor[rgb]{0.9,0.9,0.9}{
            \makecell[{{p{\linewidth}}}]{
                \texttt{\tiny{[GM$|$GM]}}
                {'word': 'преса', 'assignment': 'team'}
            }
        }
    }
    & & \\ \\

    \theutterance \stepcounter{utterance}  
    & & & \multicolumn{2}{p{0.3\linewidth}}{
        \cellcolor[rgb]{0.9,0.9,0.9}{
            \makecell[{{p{\linewidth}}}]{
                \texttt{\tiny{[GM$|$GM]}}
                {'word': 'ученик', 'assignment': 'team'}
            }
        }
    }
    & & \\ \\

    \theutterance \stepcounter{utterance}  
    & & & \multicolumn{2}{p{0.3\linewidth}}{
        \cellcolor[rgb]{0.9,0.9,0.9}{
            \makecell[{{p{\linewidth}}}]{
                \texttt{\tiny{[GM$|$GM]}}
                \texttt{opponent has won} \\
            }
        }
    }
    & & \\ \\

\end{supertabular}
}

\end{document}
